\section{Falsification Tests and Open Derivations}
\CatchFileBetweenTags{\AlphaInvVal}{constants.tex}{AlphaInvVal}
\CatchFileBetweenTags{\MeMeVPrint}{constants.tex}{MeMeVPrint}

The proposal links several Standard Model inputs to a small set of integers. Where a relation is fixed, it can fail. This section lists the sharpest tests and separates them from consequences that have been conjectured but not derived.

\subsection{Quantitative Tests}

\begin{enumerate}
    \item \textbf{The W mass.} With the measured Z mass, the on-shell mixing ratio gives $80.3845<m_W<80.3846$~GeV and an implied radiative correction $\Delta r\approx0.0351$. Standard radiative corrections on the same inputs give $m_W\approx80.356$~GeV (\cref{subsec:weak_mixing}). MeV-level W-mass measurements distinguish the two.
    \item \textbf{A fixed-input electroweak fit.} Electroweak fit packages (e.g.\ Gfitter, HEPfit) can be run with $\alpha^{-1}=\AlphaInvVal$, $\sin^2\theta_W=43/(177+16\eta)$ as the on-shell ratio, and $G_F=(\sqrt2v^2)^{-1}$ held fixed. A degradation of the global $\chi^2$ by more than $9$ relative to the floating fit would falsify the proposed fixed values. Such a fit also requires the interface response of \cref{subsec:weak_mixing}, since the proposal is not compatible with standard radiative corrections at the same inputs.
    \item \textbf{The fine-structure constant.} The proposed $\alpha^{-1}$ is a fixed number (Paper~I). Its comparison with future determinations is a direct test, and the derived scales inherit any shift.
    \item \textbf{The gravitational coupling.} $\alpha_G=(\nu-\chi/(\sigma-\chi)-\alpha)\,\alpha^{\Delta/2}$ agrees with $Gm_e^2/\hbar c$ to a few parts per million. Improved determinations of $G$ test it at that level; moving the depth by half a step changes $\alpha_G$ by a factor near twelve.
    \item \textbf{Heavy chiral partners.} On the two-wall Hamiltonian, the chiral zero modes do not take the inter-wall link $Y$ as their mass. Integrating out the gapped modes gives the light Dirac mass $M=\delta^2Y^{-1}$, so $MM^\dagger=\delta^4(Y^\dagger Y)^{-1}$ and the light masses are $\delta^2/\sigma_i(Y)$ (Paper~II, \texttt{WallSeesaw.effective\_mass}, \texttt{seesaw}). For one flavour the spectrum is exact: $E^4-(y^2+2\delta^2)E^2+\delta^4=0$, a light mode and a heavy partner with $m_{\rm light}\,m_{\rm heavy}=\delta^2$ and $m_{\rm light}^2+m_{\rm heavy}^2=y^2+2\delta^2$ (\texttt{one\_flavour\_eigen}, \texttt{light\_heavy}). Since the on-site term $\delta$ is flavour-universal, every light chiral fermion has a heavy partner with $m_{\rm heavy}=\delta^2/m_{\rm light}$: the lightest fermions have the heaviest partners, and the partner masses scale as the inverse of the observed masses ($m_{\rm heavy}(e)/m_{\rm heavy}(t)=m_t/m_e$). Observing such partners, with this inverse scaling, would test the mechanism; a light fermion without a partner at $\delta^2/m$ would contradict it. The value of $\delta$, and the identification of these zero modes with the observed fermions, are what remain to be fixed. The measured Higgs couplings constrain that identification already: they require $m\propto v$, so $\delta^2/\sigma_i(Y)$ must be linear in the VEV. If only the on-site term carried the VEV ($\delta\propto v$), every Higgs--fermion coupling would be twice its Standard Model value; if only the link did ($Y\propto v$), the couplings would have the opposite sign. LHC measurements disfavour both, so the VEV must enter $\delta$ and $Y$ together, in a way still to be derived.
\end{enumerate}

\subsection{Conjectured Consequences Requiring Derivation}

The following are consequences the proposal suggests. None is yet a theorem of the formalized framework; each needs a representation-theoretic or dynamical proof before it can serve as a prohibition.

\begin{enumerate}
    \item \textbf{No extra spatial dimensions ($D=4$).} The coordinate sector of the forced $E_8$ has rank four. Whether this excludes Kaluza-Klein modes depends on how the effective spacetime is derived from the substrate, which is open.
    \item \textbf{No supersymmetric partners ($\nu=16$).} The packet supplies the map from channels to representations: tagged by colour dimension, weak dimension and hypercharge, the sixteen channels are exactly one generation including the gauge-singlet slot (Paper~V's sterile vacancy), $(3,2)_{1/6}$, $(3,1)_{2/3}$, $(3,1)_{-1/3}$, $(1,2)_{-1/2}$, $(1,1)_{-1}$, $(1,1)_0$ (\texttt{ChannelOccupancy.occupancy}, \texttt{charge\_entries}). Every channel is used, so a further state with the same gauge numbers, such as a superpartner, has no channel in the packet (\texttt{no\_spare\_channel}). This assumes one generation occupies one packet; it does not exclude partners carried by other structures.
    \item \textbf{No exact gauge-coupling unification.} The couplings arise from distinct geometric features. This is an expectation, not a prohibition; a unified embedding has not been excluded.
    \item \textbf{Strong CP.} The package now selects $\bar\theta=0$ conditionally (\cref{sec:strong_cp_prohibition}): $\theta=0$ is the unique least-cost vacuum under $\operatorname{Re}Z$ or $|Z|$ (\texttt{PersistenceClosures.strict\_vafa\_witten}, \texttt{VacuumModulus.strict\_modulus}), and invariance under the wall-parity exchange makes the link masses Hermitian, so $\bar\theta=\theta$ (\texttt{WallParityExchange.thetaBar\_from\_exchange}). The premises are that symmetry of the update and the zero-winding sector's positive weight. A nonzero neutron EDM from strong CP, or an axion required to cancel $\bar\theta$, would contradict the selection.
    \item \textbf{No intermediate-range forces.} The proposal places force origins at the boundary and at the centroid only. Excluding other origins needs the propagation law of \cref{sec:geometric_depth}, which is not yet derived.
\end{enumerate}

\subsection{Open Derivations}

\begin{enumerate}
    \item \textbf{Workload normalization.} The actual packet trace indices are $2,2,10/3$ for weak, colour and $Y/2$. Treating the current ledgers as raw refresh weights under one uniform packet readout gives $s=5547/37777$, rather than the proposed $1849/8295$. The carrier-orbit classification now reduces symmetric incoherent preparations to six per-carrier weights; the relation needed for the proposed ratio is explicit, but the existing gates do not select those weights. Stationary diagonal preparations with nonzero up/down inter-sector hops fail that relation, while type-dependent repair lengths require a separately justified primitive task. Whether the ledgers instead already include those indices must follow from the operational work/readout map; a common calibration cannot settle the distinction (\cref{app:gauge_admittance}).
    \item \textbf{Response coefficients.} The beta-function coefficients $11$ and $2/3$, and the second-order running needed at the tau mass, should arise from the substrate's transfer matrix rather than by identification. The additive stiffness count is a stated hypothesis.
    \item \textbf{The electroweak interface response.} The factor $k$, equivalently $\Delta r\approx0.0351$, must follow from the substrate. Pure Z-pole screening overshoots it, so an anti-screening contribution is required.
    \item \textbf{The overhead's operational premises.} Register granularity, register width, the independent reference and the workload law behind $16\eta$.
    \item \textbf{The reflection charge rate.} The rate $c$ at which the scalar interface pays for reflection, which fixes the small difference between the persistent gain and $128/989$.
\end{enumerate}
